\documentclass[11pt]{article}

\usepackage{maiacustom}

\begin{document}

\psettitle{Astronomy Question Bank}

These questions were carefully written or selected to prepare students for the astronomy team selection process in Brazil. Some questions are not original; those are marked with () before the statement. The question-bank template is the same one used by Professor Kevin Zhou. His work is invaluable, and many ideas in this set can be found in his handouts.

% \begin{psidea}{Idea title}{}
% Idea
% \end{psidea}

% \begin{psexample}{Example title}{}
% Example
% \end{psexample}

% \begin{pssolution*}{}{}
% Solution
% \end{pssolution*}

% \begin{psremark*}{Remark title}{}
% Remark
% \end{psremark*}

\section{Thermodynamics}
\pts{3}
\begin{pproblem}
    A galaxy contains on the order of \(10^{10}\) stars. Because of this immense number, we can model a galaxy as a cloud of ideal gas, in which each star is equivalent to a particle of the gas.
    
    The goal of this question is to use this theoretical model to study some properties of galaxies. To that end, we make the following assumptions:

    \begin{enumerate}[label=\roman*)]
        \item The galaxy is spherical and in hydrostatic equilibrium.
        \item The mass density of the galaxy is constant and equal to \(\rho\).
        \item The stellar masses are small enough that interstellar interactions can be neglected.
    \end{enumerate}
    
    \begin{alternativas}
        \item Considering an ideal-gas system, find an expression for the pressure in terms of the density \(\rho\), the temperature \(T\), the mass of each particle \(\mu\), and physical constants.
        
        \item In our theoretical model it does not make sense to think in terms of temperature, so we need a substitute for it. Using the equipartition theorem, find an expression for \(T(r)\) and \(P(r)\).
    \end{alternativas}

\end{pproblem}

\pts{2}
\begin{pproblem}
    In this question we study the adiabatic index of stars. Assume that the exterior of a star is a vacuum at temperature \(T=0\).
    \begin{alternativas}
        \item All stars are bodies in hydrostatic equilibrium. Given that, what is the pressure at the surface of a star of mass \(M\) and radius \(R\)?
        \item Now suppose that the star expands by \(\delta R\). How does the pressure change? If needed, use \((1+x)^n\approx 1+nx\).
        \item Now conclude what minimum value of \(\gamma\), \(\gamma_{min}\), the star must have in order to remain gravitationally bound (assume that it expands adiabatically).
    \end{alternativas}


\end{pproblem}

\pts{4}
\begin{pproblem}
    There are several models for the atmosphere of our planet. Let us explore them and find the physical effects of each one.
    \begin{alternativas}
        \item First, consider the isothermal model (\(T=\) const). Given that each air particle has mass \(\mu\) and that the atmosphere has \(P(0) = P_0\), find a formula for the pressure as a function of height, \(P(h)\).
        \item A more realistic model of the atmosphere is in fact adiabatic, since air is a very poor conductor of heat. Given that the air has Poisson coefficient \(\gamma\), find an expression for \(P(h)\) in the adiabatic model. Assume that at sea level the pressure and the temperature are \(P(0)\) and \(T(0)\).
        \item Find an expression for \(dT/dh\) for the previous model and estimate its value. Is the result consistent with reality?
    \end{alternativas}


\end{pproblem}

\pts{5}
\begin{pproblem}
    Black holes are among the most fascinating bodies in the universe. In this question we study some of the thermodynamics associated with these objects. For this question we use natural units, i.e.\ \(c = G = \hbar = k_B = 1\). In this convention, the mass of the black hole is described by the equation

    \[dM = \frac{\kappa}{8\pi}dA + \Omega dL + \Phi dQ\]

    Here the quantities refer to the event horizon: \(\kappa\) is the gravitational acceleration at the event horizon, \(A\) its area, \(\Omega\) the angular velocity of the black hole, \(L\) its moment of inertia, \(\Phi\) the electric potential, and \(Q\) its charge.

    \begin{center}
    \textbf{Part 1: Basic Thermodynamics}   
    \end{center}
    
    \begin{alternativas}
        \item One of the fundamental concepts of thermodynamics is entropy. Using what you know about it, briefly explain the inequality
        \[\oint dS \geq 0\]

        \item Of the 3 factors that govern the mass of a black hole (\(dA, \ dL, \ dQ\)), only \(dA\) obeys the same inequality as the entropy. Why is this always true?
        
        For the next items, consider a black hole with no spin and no charge.

        \item Bekenstein and Hawking proved the so-called \textit{Bekenstein-Hawking entropy}, which relates the entropy to the area of the black hole (remember that we are using natural units, so some dimensions may not make sense). Bekenstein and Hawking found that for a black hole \(S \equiv \frac{A}{4}\). The equation we then have is
        
        \[dM = \frac{\kappa}{2\pi}d\left(\frac{A}{4}\right)\]

        By analogy between \(dM\) and some state function, find the temperature of the black hole as a function of \(\kappa\).

        An important term is still missing from the previous formula: the pressure associated with the dark-energy density, \(\Lambda\).

        \item The pressure due to dark energy is given by
        \[P = -\frac{\Lambda}{8\pi}\]

        where \(\Lambda\) is constant. This shows that \(VdP\) vanishes, so it may be added freely to the previous expression. We can therefore conclude that the mass of the black hole is in fact related to another thermodynamic potential. Which one is it?

        \item Note that the volume, \(V=V\), and the entropy, \(S = A/4\), are no longer independent for black holes. Assuming that the event horizon of the black hole is a sphere, find a relation between \(S\) and \(V\). This is further evidence that the internal energy, \(U = U(S, V)\), is not the best thermodynamic potential to work with.
    \end{alternativas}

    \centering
    \textbf{Part 2: Carnot Cycle for Black Holes}
    
    \begin{alternativas}
        \item The goal of this part is to build a theoretical model for a Carnot cycle inside black holes. First, however, prove an important result: for black holes, adiabatic and isochoric processes must be equivalent.
        
        \item Calculate the heat capacity at constant pressure of a black hole. Your result should be something bizarre.
        
        \item Use the fact that \(Q = T \Delta S\) along the isotherms, together with the results of the previous parts, to calculate the efficiency of a black-hole Carnot engine and confirm that you
        obtain the Carnot efficiency. Marvel at how much faster this calculation is than the
        typical derivation of the Carnot efficiency, and note that you have also inadvertently
        calculated the efficiency of the Stirling cycle.

        If you are interested in the subject, there is an interesting paper that discusses cycles in black holes specifically, available \href{https://arxiv.org/pdf/1404.5982}{here}.
    \end{alternativas}

    \centering
    \textbf{Part 3: Lifetime and Evaporation of Black Holes}
    
    \begin{alternativas}
        \item As calculated in part 1, black holes have a temperature. As a result, they emit radiation as described by the Stefan-Boltzmann law. Given that, find a relation between the lifetime of a black hole and its mass \(M\).
    \end{alternativas}

    


\end{pproblem}

\pts{3}
\begin{pproblem}
    Consider a rocket that uses an ideal diatomic gas as fuel. Its mechanism is quite simple: the gas starts in a chamber at temperature \(T_1\) with cross-sectional area \(A_1\); the gas then flows adiabatically and is expelled through an opening of area \(A_2\), at pressure \(P_2\) and temperature \(T_2<T_1\). Assuming that the flow is continuous, determine the thrust felt by the rocket.


\end{pproblem}


\pts{5} 
\begin{pproblem} (Adapted from Iran Problem Set)
    This problem aims to calculate the boiling point of liquids.

Particles in a liquid move at different speeds depending on the temperature, and some of these particles can escape the intermolecular forces and Earth's gravity (which will be neglected in this problem), leaving the surface of the liquid. These particles transfer their momentum, creating pressure when they collide with the surrounding environment. This pressure is known as the vapor pressure of the liquid. The boiling point is the temperature at which the vapor pressure equals the atmospheric pressure around the liquid.

\begin{alternativas}
\item Using the Maxwell-Boltzmann distribution, find a relation for the root-mean-square speed $v_{\text{rms}}$.

The Maxwell-Boltzmann distribution is
\begin{equation}
    n(v) \, dv = n \left(\frac{m}{2\pi k_B T}\right)^{3/2} e^{-\frac{mv^2}{2k_BT}} \, 4\pi v^2 dv
\end{equation}

The root-mean-square speed is given by
\begin{equation}
    v_{\text{rms}}^2 = \frac{1}{N} \sum_{i=1}^N v_i^2 = \frac{1}{n} \int_0^\infty v^2 n(v) dv
\end{equation}

\item Suppose that the speed of the particle under study equals $v_{\text{rms}}$. Also suppose that Earth's atmosphere is 80\% nitrogen and 20\% oxygen, and that the liquid under study is water.

Calculate the distance the escaped particle travels in the atmosphere before colliding, known as the mean free path. Express this distance in terms of the number density of the atmosphere and the geometric cross section of the particles.

\item Calculate the rate of change of the momentum of an escaping particle. Divide the change in momentum by the time interval of the process. Assume that the particles of the liquid lose all of their momentum when they collide with air molecules.

The mean time until the next collision is the mean free path divided by the speed of the particle.

Given that, over this time interval, a momentum equal to the momentum of the liquid particle is transferred to the air molecule, use Newton's second law to calculate the force exerted by the liquid particle on the air molecule.

\item Pressure is the force exerted on a surface. Find an expression for the vapor pressure of a liquid. You may leave your answer in terms of the cross sections of water and of air, \(S_w\) and \(S_a\), respectively.

\item Using the hydrostatic-equilibrium relation and assuming constant gravitational acceleration, constant air density, and zero pressure in the upper layers of the atmosphere, find a relation for the pressure near Earth's surface. Express this relation in terms of the air density, the gravitational acceleration, and the thickness of the atmosphere.

\item Impose the condition required for boiling by equating the atmospheric pressure near Earth's surface (obtained above) to the vapor pressure. Simplify the result until you obtain
\begin{equation}
    T = \frac{mhg S_w}{3k_B S_a}
\end{equation}

where $m$ is the mean mass of the air molecules, $g$ is the gravitational acceleration, $h$ is the thickness of the atmosphere, and the other parameters were described in the previous parts.

\item Determine the mean mass of the air molecules for the composition mentioned at the beginning of the problem.

\item Use the Bohr radius to estimate the ratio of the cross sections. Suppose an electron in a circular orbit around a proton, with the Coulomb force as the dominant force. Using Niels Bohr's assumption $L = n\hbar$, find the distance from the electron to the nucleus in terms of physical constants, $n$, and the atomic number $Z$.

\item Calculate the cross section of the air atoms and of the liquid atoms. Assume that each liquid molecule consists of two hydrogen atoms and one oxygen atom, and that the air particles consist of two oxygen atoms and two nitrogen atoms. Find the ratio of the cross sections $S_i/S_a$.

\item Take $g = 9.8 \, \text{m/s}^2$ and the height of the atmosphere to be $100 \, \text{km}$. Determine the boiling point of water.

\end{alternativas}

\end{pproblem}

\pts{3}
\begin{pproblem} (Adapted from Iran Problem Set)
    Dr. Shahram Abbassi is one of the most renowned Iranian scientists in the field of accretion disks. In one of his recent studies of the giant molecular cloud B32, he found a Sun-like star at the center of that cloud. According to his research, the cloud has a mass of $10^6 M_\odot$, a radius of $30 \, \text{pc}$, and a very high viscosity, so high that if the cloud were to collapse, the whole system would collapse with spherical symmetry.

    What matters most is to calculate the specific heat at constant volume ($C_v$) of this cloud. Using the data given and reasonable approximations, determine a bound on $C_v$ such that accretion is possible. Do these values depend on the mass and the radius of the cloud? What can we conclude from the result?

\end{pproblem}


\end{document}
